% English Abstract

With the continued evolution of cloud computing and microservice architectures, software operations fault diagnosis increasingly relies on comprehensive interpretation of runtime data. Logs provide semantic clues about abnormal events, while monitoring metrics describe the evolution of system states. Although prior studies have discussed log analysis, metric monitoring, knowledge augmentation, and structured reasoning separately, relatively focused validation is still lacking on how logs and monitoring metrics should be organized within a large-model-based diagnostic framework. In addition, without task-oriented data organization, large language models in software operations diagnosis may still produce outputs with insufficient evidence and unstable conclusions.

To address these issues, this thesis proposes a software operations fault diagnosis method based on large language models and software runtime logs. The method is centered on an LLM-based diagnostic framework, with log analysis serving as the primary pipeline and monitoring metrics used as auxiliary evidence for trend identification and cross-checking. Logs are organized through time-window filtering, contextual clue aggregation, and anomaly summarization, while monitoring metrics are processed through anomaly identification, trend abstraction, and textual representation. The related information is then unified into a shared analysis context. On this basis, retrieval augmentation is introduced only for knowledge-dependent scenarios, whereas structured reasoning is used to constrain analytical steps and result presentation.

For validation, this thesis builds a research-oriented evaluation setting composed of controlled samples, public samples, and supplementary simulated samples. The main experiment includes 21 controlled samples from 7 fault types and compares five progressively enhanced settings under the same Gemini 2.5 Flash model condition: metric-only input, log-only input, collaborative log-metric input, collaborative input with retrieval augmentation, and collaborative input with retrieval augmentation and structured reasoning. Supplementary experiments on public log samples and simulated production-like cases are further used to examine structured output, scenario generalization, and propagation-fault localization.

The results show that, under the controlled setting of this thesis, a log-centered input organization supported by monitoring metrics is the primary source of diagnostic improvement, with the collaborative input setting achieving a Top-1 root-cause accuracy of 76.2\%. When retrieval augmentation is added on top of collaborative input, Top-1 accuracy further rises to 85.7\%, and the gain is mainly concentrated in middleware- and database-related scenarios. Structured reasoning does not further improve Top-1 accuracy, but raises output compliance from 66.7\% to 95.2\% and presents key diagnostic evidence more clearly. These findings support the method formulation of a primary pipeline complemented by auxiliary mechanisms. It should be noted that the findings are obtained under controlled experimental conditions and should not be directly extrapolated to large-scale production environments without further validation, especially for high-noise logs, concurrent dependencies, unknown fault patterns, and complex fault propagation scenarios.
